\documentclass[11pt]{article}

\usepackage[margin=1in]{geometry}
\usepackage{amsmath,amssymb}
\usepackage{booktabs}
\usepackage{hyperref}

\setlength{\parindent}{0pt}
\setlength{\parskip}{0.55em}

\newcommand{\GB}{G^B}
\newcommand{\GS}{G^S}
\newcommand{\dG}{\Delta G}
\newcommand{\RB}{R^B}
\newcommand{\RS}{R^S}
\newcommand{\dR}{\Delta R}
\newcommand{\dT}{\Delta T}
\newcommand{\dd}{\delta}

\title{Capital Gains Realization Model\\A Minimal Specification}
\author{}
\date{May 2026}

\begin{document}
\maketitle

\section{Purpose}

The model tracks how a capital-gains reform changes the stock of unrealized
gains over time, and how that changed stock feeds back into realizations and
revenue.

The central object is not the total stock of unrealized gains. The central
object is the reform-minus-baseline difference:
\[
\dG_{a,t,k} = \GS_{a,t,k} - \GB_{a,t,k}.
\]

Here $a$ is age, $t$ is year, and $k$ is asset class. The asset classes are
equities, pass-throughs, primary homes, other homes, and real estate funds.

This setup is useful because the simulator is a repeated cross-section. It
does not literally follow the same taxpayer over time. The model instead
uses age-by-year cells as synthetic cohorts. A person aged $a$ in year $t$
who survives becomes part of age $a+1$ in year $t+1$. If the person dies,
the model either forgives, transfers, or realizes the embedded gain depending
on the policy regime.

The intuition is a bathtub. Unrealized gains sit in the tub. Realizations
drain the tub. Death either drains the tub, routes the water to heirs, or
turns it into immediate tax revenue. A tax reform changes both the drain rate
and the death treatment, so the stock begins to diverge from baseline.

\section{Baseline Objects}

Each cell $(a,t,k)$ starts with baseline quantities from the simulator:

\begin{itemize}
  \item $\GB_{a,t,k}$: baseline unrealized taxable gains.
  \item $r^B_{a,t,k}$: baseline realization rate.
  \item $\RB_{a,t,k}$: observed baseline realized gains.
  \item $m_{a,t}$: mortality rate for the tax unit cell.
  \item $\tau_{a,t}$: capital-gains marginal tax rate for the cell.
\end{itemize}

The baseline realization rate is
\[
r^B_{a,t,k} = \frac{\RB_{a,t,k}}{\GB_{a,t,k}},
\]
with the usual sparse-cell guard when $\GB_{a,t,k}=0$.

Two allocation matrices move gains across age cells:
\[
\sum_h \omega_{a \to h,t} = 1
\quad \text{for every decedent age } a.
\]

$A_{a \to h,t}$ moves survivors from age $a$ to destination age $h$.
$\omega_{a \to h,t}$ routes inherited carryover gains from decedents at age
$a$ to heirs at age $h$.

For one-year ages with a topcode at $a_{\max}$,
\[
A_{a \to h,t} =
\begin{cases}
1 & \text{if } a < a_{\max} \text{ and } h=a+1,\\
1 & \text{if } a = a_{\max} \text{ and } h=a_{\max},\\
0 & \text{otherwise.}
\end{cases}
\]

\section{Realization Behavior}

Realizations come from two distinct channels. Some sales are
weakly tax-elastic turnover: business exits, life-event-driven home sales,
divorce-induced asset divisions, forced liquidations, and portfolio
rebalancing not driven by tax timing. The holder's flexibility on these
sales is small, so the hazard governing them is approximately
policy-invariant. The remaining sales are discretionary timing decisions
that respond to the tax wedge.

The total realization rate is the sum of both channels:
\[
r_{a,t,k} = \lambda^I_{a,t,k} + r^V_{a,t,k}.
\]

$\lambda^I_{a,t,k}$ is the turnover hazard: the portion of baseline
realization treated as weakly tax-elastic and policy-invariant.
$r^V_{a,t,k}$ is the voluntary, tax-elastic rate.

$\lambda^I$ is not an alternative realization rate. It is the
policy-invariant turnover hazard used to price the value of waiting. The
model's total realization rate is still $r^S$; the split only says that
part of baseline realization is weakly tax-elastic and part is
discretionary. Both pieces are real sales that show up in revenue.

\subsection{Voluntary Realization Decision}

Normalize to one dollar of embedded gain. A representative holder in cell
$(a,t,k)$ chooses a voluntary realization fraction $r^V$:
\[
\max_{r^V \in [0,\,1-\lambda^I_{a,t,k}]} \; b_{a,k}(r^V) - P_{a,t,k} r^V.
\]

$b_{a,k}(r^V)$ is the non-tax reason to \emph{voluntarily} realize:
discretionary rebalancing, profit-taking, tax-loss harvesting on adjacent
positions. $P_{a,t,k}$ is the effective tax price of selling now instead of
waiting.

Use the marginal-benefit form
\[
b'_{a,k}(r^V) = \mu_{a,k} - \frac{1}{\eta_k}\log r^V.
\]

The first-order condition is
\[
\log r^{V*}_{a,t,k} = \eta_k(\mu_{a,k} - P_{a,t,k}).
\]

The intercept $\mu_{a,k}$ does not need to be separately estimated. It is
absorbed by anchoring the voluntary channel to its observed baseline value:
\[
r^{V,B}_{a,t,k} = r^B_{a,t,k} - \lambda^I_{a,t,k}.
\]

The unconstrained exponential response can exceed the feasible upper
bound $1-\lambda^I_{a,t,k}$ under large rate cuts. Cap the scenario
voluntary rate accordingly:
\[
r^{V,S}_{a,t,k}
=
\min\!\left\{
1-\lambda^I_{a,t,k},\;
r^{V,B}_{a,t,k}
\exp\!\left[-\eta_k\left(P^S_{a,t,k}-P^B_{a,t,k}\right)\right]
\right\}.
\]

The cap binds only when the unconstrained exponential would overshoot the
cell's available unrealized stock. The total scenario realization rate is
\[
\boxed{
r^S_{a,t,k}
=
\lambda^I_{a,t,k}
+
\min\!\left\{
1-\lambda^I_{a,t,k},\;
\left(r^B_{a,t,k}-\lambda^I_{a,t,k}\right)
\exp\!\left[-\eta_k\left(P^S_{a,t,k}-P^B_{a,t,k}\right)\right]
\right\}.
}
\]

If the reform is the baseline, $P^S=P^B$, the exponential is one, the cap
does not bind, and $r^S=r^B$ by construction.

The two limits clarify the structure. When $\lambda^I_{a,t,k}=0$, the model
collapses to a single-channel specification with full elasticity on
$r^B$ -- the previous spec. When $\lambda^I_{a,t,k}=r^B_{a,t,k}$, the
entire baseline rate is turnover and $r^S=r^B$ regardless of the tax
wedge; the cell is fully inelastic. Asset classes with high turnover
share are mechanically less responsive to rate changes, matching the
empirical pattern that housing has a lower realization elasticity than
liquid equities.

\subsection{Effective Tax Price}

Selling today costs $\tau_t$ per dollar of gain. Holding avoids that tax
today, but the dollar may still leave the holder's control later through
turnover, carryover at death, or deemed realization at death. The effective
tax price is the tax due now minus the discounted expected tax due later
through these largely policy-invariant paths:
\[
\boxed{
P_{a,t,k}(c)
=
\tau_t
-
\sum_{j=1}^{\infty} \beta^j s_{a,k,j}\tau_{t+j}
-
\sum_{j=1}^{\infty} \beta^j d_{a,j}c\tau_{t+j}.
}
\]

The terms are:

\begin{itemize}
  \item $\beta$: annual discount factor, default near $0.96$.
  \item $s_{a,k,j}$: probability of \emph{turnover} realization in future
        year $t+j$, conditional on still holding at $t$.
  \item $d_{a,j}$: probability of death in future year $t+j$ before the
        asset has otherwise exited the holder's control.
  \item $c$: share of the embedded tax burden that survives death and is
        internalized by the current holder.
\end{itemize}

Voluntary future realizations do not appear in the price. Future voluntary
choices are themselves endogenous responses to the future tax wedge, not
exogenous events with a fixed probability. Including them here would create
a fixed-point dependence between $P$ and $r^V$ that would have to be solved
iteratively. Restricting the price to turnover paths and death keeps the
bracket genuinely exogenous and avoids the recursion.

The death-regime values of $c$ are:
\[
c =
\begin{cases}
0 & \text{step-up basis},\\
\theta \in (0,1] & \text{carryover basis},\\
1 & \text{deemed realization at death}.
\end{cases}
\]

Under step-up, death wipes out the gain, so death has no future-tax cost to
the holder. Under deemed realization, death fully taxes the gain. Under
carryover, $\theta$ says how much the holder cares about the heir's eventual
tax bill.

If expected future rates are constant, $\tau_{t+j}=\tau_t$, the price factors
into a rate and a bracket:
\[
P_{a,t,k}(c) = \tau_t\left(1-M_{a,k}(c)\right),
\]
where
\[
M_{a,k}(c)
=
\sum_{j=1}^{\infty}\beta^j s_{a,k,j}
+
c\sum_{j=1}^{\infty}\beta^j d_{a,j}.
\]

The bracket $1-M_{a,k}(c)$ is the part of today's tax rate that really bites.
If future taxes are likely anyway, the extra cost of selling today is smaller.
That is why removing step-up can raise the level of realizations but lower
the sensitivity to tax-rate changes.

\subsection{Future Turnover and Death Probabilities}

The probabilities $s$ and $d$ come from a competing-risks process. Each year,
a holder either turns over the asset, dies, or keeps holding.

The probability that the asset is still held at the start of year $t+j$ is
\[
S_{a,k,j}
=
\prod_{i=0}^{j-1}
\left(1-\lambda^I_{a+i,t+i,k}-m_{a+i,t+i}\right),
\qquad
S_{a,k,0}=1.
\]

Then
\[
s_{a,k,j}=S_{a,k,j}\lambda^I_{a+j,t+j,k},
\qquad
d_{a,j}=S_{a,k,j}m_{a+j,t+j}.
\]

The model usually truncates the infinite sums after a long horizon because
discounting, mortality, and turnover make the tail negligible.

\subsection{Elasticity Implied by the Bracket and Voluntary Share}

Under constant-rate expectations,
\[
\frac{\partial P_{a,t,k}}{\partial \tau_t}
=
1-M_{a,k}(c).
\]

The voluntary-channel semi-elasticity is therefore
\[
\frac{\partial \log r^{V,S}_{a,t,k}}{\partial \tau_t}
=
-\eta_k\left(1-M_{a,k}(c)\right).
\]

Evaluated at baseline ($P^S=P^B$), the total-rate semi-elasticity is the
voluntary semi-elasticity scaled by the voluntary share of total
baseline realizations:
\[
\boxed{
\left.
\frac{\partial \log r^S_{a,t,k}}{\partial \tau_t}
\right|_{P^S=P^B}
=
-\eta_k\left(1-M_{a,k}(c)\right)
\frac{r^{V,B}_{a,t,k}}{r^B_{a,t,k}}.
}
\]

Away from baseline the same expression holds with the voluntary share
evaluated at the scenario, $r^{V,S}/r^S$, rather than at baseline. The
cap on $r^{V,S}$ also applies away from baseline; once the cap binds,
the local total semi-elasticity is zero.

Two channels of heterogeneity emerge from one structure. The bracket
$1-M_{a,k}(c)$ reflects regime-driven variation in tax responsiveness:
removing step-up shrinks the bracket and dampens elasticities. The voluntary
share $r^{V,B}/r^B$ reflects asset-class variation: housing and pass-throughs
have a high turnover share, so their total realization rate is mechanically
less sensitive to rate changes than equities. Both effects show up in the
same equation without additional parameters.

\section{Death Treatment}

Death is handled by three routing shares:
\[
\dd_{\mathrm{vanish}}+\dd_{\mathrm{route}}+\dd_{\mathrm{realize}}=1,
\qquad
\dd_{\mathrm{vanish}},\dd_{\mathrm{route}},\dd_{\mathrm{realize}}\in[0,1].
\]

\begin{center}
\begin{tabular}{lccc}
\toprule
Regime & $\dd_{\mathrm{vanish}}$ & $\dd_{\mathrm{route}}$ & $\dd_{\mathrm{realize}}$ \\
\midrule
Step-up basis & 1 & 0 & 0 \\
Carryover basis & 0 & 1 & 0 \\
Deemed realization & 0 & 0 & 1 \\
\bottomrule
\end{tabular}
\end{center}

Only the routed share enters heirs' future unrealized gains. The death-flow
contribution to next-period stock is
\[
\boxed{
\dG^{\mathrm{death}}_{h,t+1,k}
=
\dd_{\mathrm{route}}
\sum_a
\omega_{a \to h,t}m_{a,t}
\left(\GB_{a,t,k}+\dG_{a,t,k}\right).
}
\]

The deemed-realization share does not enter future stock. It becomes immediate
death-event revenue:
\[
\boxed{
R^{\mathrm{death}}_{t,k}
=
\dd_{\mathrm{realize}}
\sum_a
m_{a,t}
\left(\GB_{a,t,k}+\dG_{a,t,k}\right)
\overline{\tau}_{a,t}.
}
\]

$\overline{\tau}_{a,t}$ is the applicable average capital-gains rate for
deemed realizations in the cell, after any regime-specific exemptions or rate
rules.

\section{Stock Law of Motion}

The survivor channel is the main stock-flow equation. Survivors keep only the
part of their gain stock that is not realized during the year. The
reform-minus-baseline survivor contribution is
\[
\boxed{
\dG^{\mathrm{surv}}_{h,t+1,k}
=
\sum_a A_{a \to h,t}(1-m_{a,t})
\left[
(1-r^S_{a,t,k})\dG_{a,t,k}
+
\GB_{a,t,k}(r^B_{a,t,k}-r^S_{a,t,k})
\right].
}
\]

The first term inside the brackets carries forward the policy-induced stock
already present at the start of the year. The second term creates new stock
difference because the reform changes this year's realization rate. If
$r^S<r^B$, more gains stay locked in and $\dG$ rises. If $r^S>r^B$, the
reform realizes more gains and drains the stock.

Adding survivor flow and death routing gives the full recurrence:
\[
\boxed{
\begin{aligned}
\dG_{h,t+1,k}
=\;&
\sum_a A_{a \to h,t}(1-m_{a,t})
\left[
(1-r^S_{a,t,k})\dG_{a,t,k}
+
\GB_{a,t,k}(r^B_{a,t,k}-r^S_{a,t,k})
\right] \\
&+
\dd_{\mathrm{route}}
\sum_a
\omega_{a \to h,t}m_{a,t}
\left(\GB_{a,t,k}+\dG_{a,t,k}\right).
\end{aligned}
}
\]

This is the core model. Given all age cells in year $t$, it computes all age
cells in year $t+1$ in one pass.

Three checks matter:

\begin{itemize}
  \item If reform equals baseline, then $r^S=r^B$, $\dd_{\mathrm{route}}=0$,
        and $\dG$ remains zero.
  \item Under carryover, death moves unrealized gains to heirs rather than
        erasing them.
  \item Under deemed realization, death removes the stock and records revenue
        immediately.
\end{itemize}

\section{Realizations and Revenue}

The structural realization rate is applied as a ratio to observed baseline
realizations:
\[
\boxed{
\RS_{a,t,k}
=
\RB_{\mathrm{obs}}(a,t,k)
\frac{
r^S_{a,t,k}\GS_{a,t,k}
}{
r^B_{a,t,k}\GB_{a,t,k}
}.
}
\]

Since $\GS=\GB+\dG$, the change in realized gains is
\[
\boxed{
\dR_{a,t,k}
=
\RB_{\mathrm{obs}}(a,t,k)
\left[
\frac{
r^S_{a,t,k}(\GB_{a,t,k}+\dG_{a,t,k})
}{
r^B_{a,t,k}\GB_{a,t,k}
}
-1
\right].
}
\]

The ratio form is deliberate. It guarantees exact baseline reproduction when
there is no reform: when $r^S=r^B$ and $\dG=0$, the ratio is one and
$\RS=\RB_{\mathrm{obs}}$. In cells where $r^B_{a,t,k}\GB_{a,t,k}=0$, use the
standard sparse-cell fallback: either apply $r^S\GS$ directly or pool the
baseline rate within asset class.

Revenue from ongoing realizations is
\[
T^S_{a,t,k}=\tau^S_{a,t}\RS_{a,t,k},
\qquad
T^B_{a,t,k}=\tau^B_{a,t}\RB_{a,t,k}.
\]

The reform-minus-baseline revenue change is
\[
\boxed{
\dT_{a,t,k}
=
\tau^S_{a,t}\RS_{a,t,k}
-
\tau^B_{a,t}\RB_{a,t,k}.
}
\]

Total revenue for year $t$ adds ongoing realization revenue and any
deemed-realization revenue at death:
\[
\boxed{
\dT_t
=
\sum_{a,k}\dT_{a,t,k}
+
\sum_k R^{\mathrm{death}}_{t,k}.
}
\]

\section{Calibration}

The model has a small set of inputs. Most are observed or calibrated to
baseline data. Two are judgmental: the carryover bequest parameter and the
turnover share by asset class.

\begin{center}
\begin{tabular}{lll}
\toprule
Object & Role & Source \\
\midrule
$r^B_{a,t,k}$ & Baseline realization rate & Simulator baseline \\
$m_{a,t}$ & Mortality & Cell data or life table \\
$\eta_k$ & Voluntary tax-price response & Calibrated by simulation \\
$\lambda^I_{a,t,k}$ & Turnover hazard & Share of $r^B$ by asset class \\
$\beta$ & Discount factor & Default near $0.96$ \\
$c$ & Death burden share in behavior & Regime rule \\
$\theta$ & Carryover bequest motive & Tunable assumption \\
$\omega_{a\to h,t}$ & Heir age routing & Estate module \\
\bottomrule
\end{tabular}
\end{center}

\subsection{Turnover Hazard}

The turnover hazard captures sales that are weakly tax-elastic: business
exits, life-event home sales, divorce-induced asset divisions, forced
liquidations, non-tax-driven portfolio rebalancing. The natural calibration
is as a share of the baseline realization rate:
\[
\lambda^I_{a,t,k} = \phi^I_k\, r^B_{a,t,k}, \qquad \phi^I_k\in[0,1).
\]

The share form is automatically non-negative, automatically respects the
ceiling $\lambda^I_{a,t,k}\le r^B_{a,t,k}$, and degrades gracefully in
sparse cells where $r^B$ is small or zero. Allowing the turnover share
$\phi^I_k$ to depend on age is a reasonable refinement but not pursued
in the default calibration.

Three calibration paths are reasonable, in increasing order of effort:

\begin{enumerate}
  \item \textbf{Default share by asset class.} Pick a single $\phi^I_k$
        and use the share form above. Reasonable starting values:
        $\phi^I_k\approx 0.3$ for liquid equities, $0.7$ for primary
        and other homes, $0.5$ for pass-throughs and real estate funds.
  \item \textbf{Implied from elasticity heterogeneity.} If the
        asset-class long-run target elasticity is $\varepsilon_k$ and a
        common voluntary structural elasticity is
        $\varepsilon^V=-\eta(1-M(0))$, the implied turnover share is
        \[
        \phi^I_k = 1 - \frac{\varepsilon_k}{\varepsilon^V}.
        \]
        This route makes cross-asset elasticity differences the explicit
        drivers of cross-asset turnover shares.
  \item \textbf{External data.} Use mortgage origination and
        life-event-driven sale rates from real-estate panels for housing,
        business-exit rates from BEA / Census Business Dynamics
        Statistics for pass-throughs, and survey data on non-tax-driven
        portfolio rebalancing for equities. Convert the observed
        absolute hazard to a share by dividing by $r^B$, capped at one.
\end{enumerate}

\subsection{Tax-Price Curvature}

The closed-form aggregate semi-elasticity expression
\[
\left.
\frac{\partial \log \sum_{a,t}\RS_{a,t,k}}{\partial \tau}
\right|_{\tau=\tau_0,\;c=0}
\;=\;
-\eta_k
\left(1-M_{a,k}(0)\right)_{\mathrm{aggregated}}
\frac{r^{V,B}_k}{r^B_k}
\]
is the right intuition for what $\eta_k$ does, but it is not a precise
calibration target. With the cap on $r^{V,S}$, sparse-cell pooling,
stock-feedback dynamics, and weighted aggregation, the actual aggregate
elasticity of $\sum\RS$ with respect to $\tau$ differs from the closed
form. Calibrate $\eta_k$ by simulation: perturb $\tau$ by a small
amount under current law, re-run the bathtub recurrence, and find the
$\eta_k$ that matches the target aggregate semi-elasticity. The
voluntary-share scaling implies the simulation-calibrated $\eta_k$ is
larger in magnitude than in a single-channel calibration that hits the
same target, because part of the baseline rate is held inelastic.

\subsection{Death Burden Share}

The default death burden share is
\[
c =
\begin{cases}
0 & \text{step-up},\\
\theta & \text{carryover},\\
1 & \text{deemed realization}.
\end{cases}
\]

A reasonable starting value for carryover is $\theta=0.5$, with sensitivity
tests around it. The interpretation is simple: the holder behaves as if half
of the heir's future capital-gains tax matters for the holder's own decision.

\section{Cell-Level Conventions}

For joint filers, assign the household to the older spouse's age:
\[
a_{\mathrm{joint}}=\max(a_1,a_2).
\]

The joint-filer mortality rate is the probability that the joint-filer unit
ceases to exist as a joint unit:
\[
m_{\mathrm{joint},t}=m_{1,t}m_{2,t}.
\]

Within an age-year-asset cell, distribute $\dR_{a,t,k}$ to records in this
order:

\begin{enumerate}
  \item Pro rata to baseline realized gains, if the cell has baseline
        realized gains.
  \item Pro rata to baseline unrealized gain stock, if baseline realizations
        are zero but gain stock is positive.
  \item Assign nothing if both baseline realizations and baseline gain stock
        are zero.
\end{enumerate}

The youngest age cell has no survivor inflow. It can receive unrealized gains
only through inheritance under carryover. The topcoded age cell loops into
itself.

\section{Model Boundaries}

The model is meant to capture realization timing, lock-in, carryover routing,
and deemed-realization revenue. It does not try to solve every capital-gains
problem.

Five limits should be kept in view:

\begin{itemize}
  \item Accruals on existing assets are treated as policy-invariant, so they
        difference out of $\dG$.
  \item The turnover / voluntary partition is a structural assumption.
        In reality the line is fuzzy: a divorce-driven home sale could be
        deferred a year if the tax cost is large enough. The model
        collapses that fuzziness into a single hazard $\lambda^I$ that
        is treated as policy-invariant.
  \item Charitable transfers of appreciated assets are not modeled as a
        separate stock drain.
  \item Announced future tax rates can enter through the path
        $\tau_{t+j}$, but unannounced future policy risk is not forecast
        endogenously.
  \item Mark-to-market regimes are outside this specification because annual
        taxation of accruals changes the revenue channel rather than just the
        realization timing channel.
\end{itemize}

\end{document}
